In recent years there has been an increased interest within the high energy and nuclear physics communities in the application of hadron colliders as photon-photon and photon-hadron colliders. The charge of the accelerated hadron generates a photon flux which allows for measurements that utilize photons as probes into QCD dynamics, nuclear structure, and saturation scale physics. For my research I propose to study these photon interactions using the ATLAS experiment at the Large Hadron Collider (LHC). In this document we will discuss photon interactions at the LHC, and the types of ultra peripheral  (UPCs) ion-ion collisions which enable them. Then we will delve into open questions regarding specific collision geometries, such as photon emitter ambiguities in events without nuclear breakup, and how to address them in di-jet events. I will propose using machine learning models to differentiate between different UPC photoproduction sub-processes and how to use that framework to extract additional information regarding these events.